%!TEX root = mainQBN.tex
\section{Introduction}

Pearl's \emph{Bayesian networks}~\cite{Pearl86,Pearl_book} provide a framework for reasoning under conditions of uncertainty and partial knowledge, with a wide range of applications from statistics to epidemiology, economics and computer science.
Bayesian networks have a \emph{dual nature}, serving both as probabilistic graphical models for classical probabilistic reasoning and inference, and as causal models, precising the connections between observed data and causal relations.
When reasoning on \emph{quantum systems}, the classical framework is not general enough to account for entanglement and the non-local correlations observed in Bell experiments.
The development of quantum causal models (see \eg~\cite{BarrettLorenzOreshkov2019} and references therein) is an active research area across quantum information and the foundations of quantum theory, advancing along various axes, whose motivations span from foundational questions and non-locality, to enabling device-independent cryptographic protocols, to facilitating data-driven discovery.

In this paper, we focus on quantum Bayesian networks, a direct generalization of Pearl's networks introduced in foundational work by Henson, Lal, and Pusey~\cite{GBN2014}.
They provide a mathematical framework to describe causal relations, to analyze correlations, and to predict the probabilities of measurement outcomes, in systems involving both classical and quantum data.
The formalism builds on previous work by Leifer and Spekkens~\cite{LeiferSpekkens2013}, where the perspective is that of \emph{quantum theory as a theory of inference}.
Quantum theory is indeed fundamentally probabilistic at its core, as it is concerned with predicting the probabilities of measurement outcomes on a physical system.
In this sense, the prediction task can be framed as a problem of \emph{probabilistic inference over models involving both classical and quantum data}, as \Cref{ex:bell} illustrates.
Probabilistic inference then offers mathematical and logical tools to comprehend, predict, and control quantum phenomena, essential both to the theoretical understanding and to quantum technologies.

\begin{example}[Alice \& Bob: the Bell set-up]
\label{ex:bell}
The directed acyclic graph in \cref{fig:Bell} describes the well-known set-up for the Bell experiment.
Alice and Bob---who stand in widely separated laboratories---are each able to perform two possible measurements on a qubit (for example, measuring it with respect to two different bases). 
Their colleague Quentin prepares a pair of (possibly entangled) qubits, sending one to Alice and the other to Bob. When Alice receives her qubit $q_1$, she chooses to randomly perform one of the two possible measurements, by flipping a coin $X$.
When Bob receives his qubit $q_2$, he also performs a measurement, by flipping a coin $Y$.
The result of the experiment is
\begin{equation}\label{eq:bell}
	\Pr(a,b\mid x,y) =\Pr(a,b,x,y) /\Pr(x,y)
\end{equation}
\ie\ the probability that the (classical) outcomes of Alice and Bob measurements are respectively $a$ and $b$, given outcomes $x$ for $X$ and $y$ for $Y$.
\end{example}

\begin{figure}
	\begin{minipage}[b]{0.39\linewidth}
		\begin{adjustbox}{}
			\begin{tikzpicture}
				\begin{genscope}
					\node (a) at (-1.5,0) {$A$};
					\node (b) at (1.5,0) {$B$};
					\node (c) at (0,1) {$\qreg$};
					\node (x) at (-3,1) {$X$};
					\node (y) at (3,1) {$Y$};
					
					\path (y) edge (b);
					\path (x) edge (a);
					\path (c) edge node[named path]{$q_2$} (b);
					\path (c) edge node[named path]{$q_1$} (a);
				\end{genscope}
			\end{tikzpicture}
		\end{adjustbox}
		\caption{Bell set-up (from~\cite{GBN2014}).}
		\label{fig:Bell}
	\end{minipage}
	\hfill
	\begin{minipage}[b]{0.59\linewidth}
		\begin{adjustbox}{}
			\begin{tikzpicture}[baseline=(base)]
				\coordinate (base) at (0,0);
				\begin{scope}[every node/.style={rectangle,draw=none}, every path/.style={draw=black}]
					\node[draw=black,minimum width=2.5em] (bQ) at (1,0) {$\qreg$};
					\node[draw=black,minimum width=2.5em] (bX) at (3,0) {$X$};
					\node[draw=black,minimum width=2.5em] (bA) at (6,0) {$A$};
					\node[draw=black,minimum width=2.5em] (bB) at (8,0) {$B$};
					\node[draw=black,minimum width=2.5em] (bY) at (10,0) {$Y$};
					
					\node (cX) at (4.5,-1.5) {$\cut$};
					\node (cQ) at (4,-2.25) {$\cut$};
					\node (pQ) at (6.7,-1.5) {$\parr$};
					\node (cY) at (9.1,-1.5) {$\cut$};
					\path (3.3,-.25) |- node[named edge,pos=.5,above right]{$X^+$} (cX) -| node[named edge,pos=.5,above left]{$X^-$} (5.7,-.25);
					\path (1.3,-.25) |- node[named edge,pos=.5,above right]{$\qreg^+$} (cQ) -| node[named edge,pos=.5,above left]{$\qreg^-$} (pQ);
					\path[out=180,in=-90] (pQ) edge node[named edge,pos=.4,right]{$q_1^-$} (5.9,-.25);
					\path[out=0,in=-90] (pQ) edge node[named edge,pos=.4,left]{$q_2^-$} (7.7,-.25);
					\path (10.3,-.25) |- node[named edge,pos=.5,above left]{$Y^+$} (cY) -| node[named edge,pos=.5,above right]{$Y^-$} (7.9,-.25);
					\path (6.3,-.25) -- node[named edge,pos=.75]{$A^+$} ++(0,-.6);
					\path (8.3,-.25) -- node[named edge,pos=.75]{$B^+$} ++(0,-.6);
				\end{scope}
			\end{tikzpicture}
		\end{adjustbox}
		\caption{Bell set-up as a quantum proof-net.}
		\label{ex:bell_pn}
	\end{minipage}
\end{figure}

\subparagraph{Bayesian Networks.}

In the theory of Bayesian networks, the causal structure is encoded by a \emph{directed acyclic graph} (DAG) where nodes represent random variables and edges express conditional dependencies.
The strength of the dependencies (or the degree of knowledge) is quantified by conditional probability tables.
A key benefit of Bayesian networks is to provide a compact representation of large probability distributions, and efficient inference algorithms (both exact and approximate) to answer queries about the underlying distribution without explicitly constructing it in full.
Critically, the DAG includes both \emph{observable} variables of interest, and \emph{hidden} (unobserved) ones; in figures, we adopt the common convention of denoting classical hidden variables by shaded nodes. In the quantum setting---as in Bell's theorems---hidden variables play a central role, and they may correspond to quantum systems.

\subparagraph{Semantics and Inference.}

The semantics of a Bayesian network is the probability distribution it defines.
More accurately, what one seeks is the \emph{marginal distribution over the variables of interest}.
Exact inference computes it precisely---this involves two key operations:

{\hfill
	product (\emph{composing}) 
	+ summing-out (\emph{hiding}) irrelevant variables.
\hfill}

The formalization and theory of inference rely on a class of functions, known as factors (see \Cref{sec:BN}), which is an abstraction of conditional probability distributions.
Tractability and efficiency rely on their properties, and specifically on two key aspects: the product of factors \emph{inherently shares variables}, and product and sum \emph{distribute} under suitable conditions---pushing the sum on smaller components reduces the size of computations.

\subparagraph{Quantum Bayesian Networks, issues.}

Quantum Bayesian networks are still an emerging field, not as developed as their classical counterparts.
A crucial missing feature is the ability to compute the semantics of the model (the desired marginal distribution) through intermediate, partial computations, without ever computing the full joint distribution.
Put differently, what is lacking is \emph{compositionality}, which would enable computing a model's denotation as a function of its subparts, alongside modular reasoning.
A closely related question concerns \emph{modularity}: when can causal descriptions of systems as subparts be used to construct larger models?
A well-established tool to ensure modularity are \emph{types}, that specify precise contracts (\eg\ input/output behaviors) for the components of a system. 

\subparagraph{The goal of this paper.}

We have two main objectives in this paper.
\begin{itemize}
\item
To address the lack of compositionality and modularity in the setting of quantum Bayesian networks by introducing methods and concepts from denotational semantics and proof theory, thereby enabling compositional principles and a typing discipline.
\item
To develop a framework fully compatible with Bayesian networks and Bayesian inference, thus paving the way for the application of the techniques developed in that context.
\end{itemize}

\emph{Compositionality and modularity} facilitate reasoning about complex systems and their properties, ensuring that the meaning of the system can be derived systematically and in a principled way from the meanings of its parts, and allowing components to be analyzed and replaced independently.
Compositionality also enables modular reasoning about inference.

\emph{Types} serve as abstract characterizations of systems behaviors, constraining and guiding their formation.
Types discipline statically guarantees semantic properties such as termination, consistency, and compositional correctness.
Well-typed programs exhibit well-behaved execution and preserve quantitative (\eg\ probabilistic) invariants throughout their evaluation.

\emph{Encompassing the semantics of Bayesian networks.}
As seen in \cref{ex:bell} (the Bell set-up), even for systems involving quantum sources of causality, only the classical outcomes of measurements can be observed.
Thus, what a model ultimately defines is a probability distribution over classical variables---such as \cref{eq:bell}.
From the inference perspective, it is desirable to have a framework enabling (when relevant) the extensive set of inference techniques and algorithms developed for Bayesian networks.
To this end, our semantics integrates a key notion of inference algorithms, that of \emph{factor}, which we discuss in \cref{sec:BN}.

\subparagraph{Contributions and challenges.}

Our first contribution is to develop a \emph{compositional semantics}, which allows for the interpretation and modular combination of components.
We adapt Selinger's semantics in~\cite{Selinger2004} to take into account the factor-based approach of Bayesian networks.
The main technical challenge is to conciliate two very different behaviors:
\begin{itemize}
\item
\emph{Classical variables share their values, and they do so in an efficient way}: the mathematical setting underlying Bayesian networks integrates this feature in the very definition of factors product, and exploits it to obtain compact representations and efficient calculations.
\item
\emph{Quantum data cannot be shared}: a defining feature of quantum computing is the \emph{No-Cloning Theorem}, implying qubits cannot be duplicated nor broadcast to multiple receivers.
\end{itemize}

We satisfy both requirements by introducing \emph{quantum factors} in \cref{sec:semantic}.
The mathematical developments in that section are our main and most technical results. 
Remarkably, when all causes are classical, our framework exactly coincides with the standard \emph{factor-based} semantics of Bayesian networks, while in the purely quantum case it behaves like tensor networks.
In \cref{sec:qbn}, we rely on quantum factors to redefine quantum Bayesian networks.
Our formalism is equivalent to that in~\cite{GBN2014}, however our semantics enables a compositional interpretation by sub-components, unlike the original definition (see~\cite[page 12]{GBN2014}), as we discuss in \cref{sec:lack}.
A crucial aspect to achieve compositionality is that quantum factors are closed under both product and sum-out (\ie\ marginalization of unobserved variables).
Finally, in \cref{sec:qpns}, we explore a \emph{typed graphical formalism} based on proof-nets of linear logic, where types ensure well-behaved compositions of systems.
Our key result is that the formalism is \emph{sound and complete} \wrt\ quantum Bayesian networks: every quantum Bayesian network can be represented as a proof-net, and every closed proof-net corresponds to a quantum Bayesian network.

\paragraph*{Motivational examples: compositionality and modularity}
\label{sec:motivational}

Two examples can illustrate the desiderata and issues with compositionality and modularity.

\subparagraph{Compositionality.}

Consider the DAG in \Cref{fig:compositional}.
The nodes $X$ and $Y$ produce a classical output, while the nodes $Q_1$ and $Q_2$ have a quantum nature.
A natural question is:
\emph{Can we compute the semantics of the model in terms of sub-components}---as for example the highlighted sub-graphs? 
The approach in~\cite{GBN2014} does not adapt well to an interpretation by components because it has a global nature, as we will discuss in \Cref{sec:lack}.

\subparagraph{Do parts compose well?}

When considering components (possibly with open inputs\footnote{Bayesian networks with open inputs are called \emph{conditional Bayesian networks}, see \eg~\cite{KollerBook}.}), a natural question somehow dual to the previous one is whether independently defined components compose well.
Consider the three DAGs in \Cref{fig:modular}: $\N_0$ awaits an input $A$ and outputs $C$, $\N_1$ and $\N_2$ both await an input $C$ and output $A$ and $D$.
The graph obtained by plugging together $\N_0$ and $\N_1$ (matching inputs and outputs) is a DAG, while the graph that plugs together $\N_0$ and $\N_2$ has a directed cycle.
By moving to \emph{typed graphs}, we guarantee that composing graphs of compatible types produces a DAG (\cref{sec:PNcompositionality}).
Rather than defining yet-another-syntax, we encode quantum Bayesian networks into the graph syntax of linear logic: proof-nets.
This builds on a recent line of work connecting (classical) Bayesian networks with proof-nets~\cite{EhrhardFP23,popl24,BPN}.
A \pn is typed by a sequent in Multiplicative Linear Logic.
As we will see (\cref{sec:PNcompositionality} and \cref{fig:modularity_pn}), the DAGs in \cref{fig:modular} admit the following typing:
\begin{equation*}
\hfill \N_0\der A\lollipop C \hfill \N_1\der (A\lollipop C)\lollipop D \hfill \N_2\der C \lollipop (A \otimes D) \hfill
\end{equation*}
The DAGs $\N_0$ and $\N_1$ compose together, producing a DAG of output $D$, while $\N_2$ cannot be given any type that matches the one of $\N_0$.

\begin{figure}
\begin{minipage}[b]{0.35\linewidth}
\begin{adjustbox}{max height = 8em, center = \linewidth}
\begin{tikzpicture}[bg_ellipse/.style={ellipse,minimum height=1.8cm,inner sep=0pt}]
		\coordinate (T_Q1) at (0, 0);
		\coordinate (T_Y) at (2.5, 0);
		\coordinate (T_Q2) at (1.5, -1.5);
		\coordinate (T_X) at (-0.5, -3.0);
		
		\node[bg_ellipse, fill=red!10, rotate=80.5, minimum width=4.5cm] at ($(T_Q1)!0.5!(T_X)$) {};
		\node[bg_ellipse, fill=blue!10, rotate=56.3, minimum width=3.2cm] at ($(T_Y)!0.5!(T_Q2)$) {};
\begin{genscope}
		\node (Q1) at (T_Q1) {$\qreg_1$};
		\node[fill=gray!50] (Y) at (T_Y) {$Y$};
		\node (Q2) at (T_Q2) {$\qreg_2$};
		\node (X) at (T_X) {$X$};
		
		\path (Q1) edge (X);
		\path (Q1) edge (Q2);
		\path (Y) edge (Q2);
		\path (Q2) edge (X);
		\path (X) edge ++(0, -0.8);
\end{genscope}
\end{tikzpicture}
\end{adjustbox}
\caption{Compositionality.}
\label{fig:compositional}\label{fig:components}
\end{minipage}
\hfill
\begin{minipage}[b]{0.6\linewidth}
\begin{adjustbox}{max height = 8em, center = \linewidth}
\begin{tikzpicture}[baseline=(base)]
	\node (base) at (-1, 2.5) {$\N_0$};
\begin{genscope}
	\node[fill=gray!50] (B) at (0, 2) {$B$};
	\node (C) at (0, .5) {$C$};
	
	\path (0, 3.25) edge node[named edge,right]{$A$} (B);
	\path (B) edge (C);
	\path (C) edge node[named edge,right]{$C$} (0, -.5);
\end{genscope}
\end{tikzpicture}
\quad\vrule\quad
\begin{tikzpicture}[baseline=(base)]
	\node (base) at (-1.5, 3) {$\N_1$};
\begin{genscope}
	\node[fill=gray!50] (E) at (0, 3) {$E$};
	\node (D) at (-1.5, 1) {$D$};
	\node (A) at (1.5, 2) {$A$};
	
	\path (E) edge (D);
	\path (E) edge (A);
	\path ($(D)+(-0.8, 1.0)$) edge node[named edge,left]{$C$} (D);
	\path (D) edge node[named edge,left]{$D$} ($(D)+(-0.6, -1.0)$);
	\path (A) edge node[named edge,right]{$A$} ($(A)+(0.5, -1.5)$);
\end{genscope}
\end{tikzpicture}
\quad\vrule\quad
\begin{tikzpicture}[baseline=(base)]
	\node (base) at (-1.5, 3) {$\N_2$};
\begin{genscope}
	\node[fill=gray!50] (E) at (0, 3) {$E$};
	\node (D) at (-1.5, 1) {$D$};
	\node (A) at (1.5, 2) {$A$};
	
	\path (E) edge (D);
	\path (E) edge (A);
	\path ($(D)+(-0.8, 1.0)$) edge node[named edge,left]{$C$} (D);
	\path (D) edge node[named edge,left]{$D$} ($(D)+(-0.6, -1.0)$);
	\path (A) edge node[named edge,right]{$A$} ($(A)+(0.5, -1.5)$);
	\path[out=-50, in=-100] (D) edge (A);
\end{genscope}
\end{tikzpicture}
\end{adjustbox}
\caption{Modularity.}
\label{fig:modular}
\end{minipage}
\end{figure}
